\section{Two Questions, and Only One Answer Each}

Here is a sentence from the \emph{Catechism of the Catholic Church}, printed without visible strain: ``In a sense bodily death is natural, but for faith it is in fact `the wages of sin.'\,''\endnote{CCC 1006, Holy See English web text, read 2026-07-26. The paragraph opens by quoting \emph{Gaudium et spes} 18 (``It is in regard to death that man's condition is most shrouded in doubt'') and closes: ``For those who die in Christ's grace it is a participation in the death of the Lord, so that they can also share his Resurrection.'' The footnoted authority for ``the wages of sin'' is Romans 6:23 with Genesis 2:17.} Two paragraphs later: ``Even though man's nature is mortal God had destined him not to die.''\endnote{CCC 1008, read as above. The paragraph's full sequence: ``Death is a consequence of sin. the Church's Magisterium, as authentic interpreter of the affirmations of Scripture and Tradition, teaches that death entered the world on account of man's sin. Even though man's nature is mortal God had destined him not to die. Death was therefore contrary to the plans of God the Creator and entered the world as a consequence of sin.'' The quoted phrase in the last sentence---``bodily death, from which man would have been immune had he not sinned''---is from \emph{Gaudium et spes} 18, and the ``last enemy'' is 1 Corinthians 15:26. The Catechism's own delivery lowercases some sentence-initial articles (``the Church's Magisterium''); quotations here reproduce its wording and are begun so that the casing of the quoted words is unaffected.}

Read those two flatly and they are hard to reconcile. Man's nature is mortal; death entered the world through sin; death was contrary to God's plan. If a nature is mortal, death is what it does, and something a nature does is not obviously an intruder. Either the Catechism is being sloppy, or it is holding two claims apart with a distinction it has not troubled to spell out. It is doing the second, the distinction is very old, and unpacking it is the single most useful piece of work this essay can do---because almost every popular confusion about Christian teaching on death comes from collapsing it.

\subsection{The texts, read where they stand}

Start with the loci themselves rather than with a chain of quotations arranged to prove something.

Genesis puts the threat in the mouth of God in the garden: ``But of the tree of knowledge of good and evil, thou shalt not eat. For in what day soever thou shalt eat of it, thou shalt die the death.''\endnote{Genesis 2:17, Douay--Rheims, read 2026-07-26 in the tracked Challoner verse text. ``Thou shalt die the death'' renders the Latin \latin{morte morieris}, itself a rendering of the Hebrew infinitive absolute construction; the doubling is emphatic in the source languages and reads as a Hebraism in English. This essay does not examine the Hebrew and makes no claim about the idiom's force beyond noting that the Latin and English preserve a doubling that is not native to either.} The sentence, when it comes, is a different kind of statement: ``In the sweat of thy face shalt thou eat bread till thou return to the earth out of which thou wast taken: for dust thou art, and into dust thou shalt return.''\endnote{Genesis 3:19, Douay--Rheims, read as above. The observation in the text---that the sentence is stated as a fact about the material the man is made of---is a reading of the verse's own grammar (\latin{quia pulvis es}, ``for dust thou art''), offered as such.} Notice the form of that. The penalty is announced in a clause that gives its own reason, and the reason is not \emph{because you sinned} but \emph{because you are dust}. The verse imposes a sentence by pointing at a material fact. Whatever else Genesis is doing, it is not treating the mortality of a body made from earth as a novelty introduced by the transgression; it is treating the transgression as the reason that fact is now going to be allowed to operate.

The book of Wisdom takes the other side of the tension with equal firmness: ``For God made not death, neither hath he pleasure in the destruction of the living\ldots{} For God created man incorruptible, and to the image of his own likeness he made him. But by the envy of the devil, death came into the world.''\endnote{Wisdom 1:13 and 2:23--24, Douay--Rheims, read 2026-07-26 in the tracked Challoner verse text. Wisdom is a late book written in Greek, and the \emph{death} it has chiefly in view in these chapters is arguably the spiritual death of the wicked, given 1:16 (``But the wicked with works and words have called it to them\ldots{} because they are worthy to be of the part thereof'') and 2:24's conclusion in the same verse: ``and they follow him that are of his side.'' This essay records the ambiguity rather than resolving it. The verse's dogmatic career, however, does not depend on the resolution: Aquinas and the Catechism both take \latin{Deus mortem non fecit} as bearing on bodily death, and the essay follows their use while flagging that the exegesis is contestable.} And Paul supplies the formula the whole later tradition works from: ``Wherefore as by one man sin entered into this world and by sin death: and so death passed upon all men, in whom all have sinned''; and, more compactly, ``For the wages of sin is death.''\endnote{Romans 5:12 and 6:23, Douay--Rheims, read 2026-07-26 in the tracked Challoner verse text. The Douay's ``in whom all have sinned'' follows the Vulgate's \latin{in quo omnes peccaverunt}, and it is that Latin form which the Council of Trent quotes in its decree on original sin. The construction underlying the Latin relative is disputed among interpreters; this essay did not examine the Greek and therefore asserts nothing about it, relying on Romans 5:12 only for the association of sin's entry with death's entry, which is not in dispute on any reading of the clause.}

Two scriptural facts, then, not one. The tie between sin and death is asserted in the strongest terms available. And the mortality of a creature made of dust is asserted as the ground of the sentence rather than as its content. Any account that suppresses either half is suppressing Scripture.

\subsection{What Trent defined, and what it did not}

The Church's most exact statement is the first canon of the fifth session of the Council of Trent, promulgated 17 June 1546, and it repays being read word by word rather than summarized.

``If anyone does not confess that the first man, Adam, when he had transgressed the commandment of God in Paradise, immediately lost the holiness and justice wherein he had been constituted; and that he incurred, through the offence of that prevarication, the wrath and indignation of God, and consequently death, with which God had previously threatened him, and, together with death, captivity under his power, who thenceforth had the empire of death, that is to say, the Devil; and that the entire Adam, through that offence of prevarication, was changed, in body and soul, for the worse; let him be anathema.''\endnote{Council of Trent, session V (17 June 1546), decree concerning original sin, canon 1; English from the public-domain translation of J.~Waterworth (London: Burns and Oates; reprint of the London 1848 translation), and the Latin from the Tauchnitz stereotype reprint of the Roman edition (Leipzig, 1887). Both were read 2026-07-26 in the full texts of those printings as digitized by the Internet Archive at the items registered for those editions in the source library. The Latin: ``\latin{Si quis non confitetur, primum hominem Adam, quum mandatum Dei in paradiso fuisset transgressus, statim sanctitatem et iustitiam, in qua constitutus fuerat, amisisse, incurrisseque per offensam praevaricationis huiusmodi iram et indignationem Dei, atque ideo mortem, quam antea illi comminatus fuerat Deus, et cum morte captivitatem sub eius potestate, qui mortis deinde habuit imperium, hoc est diaboli, totumque Adam per illam praevaricationis offensam secundum corpus et animam in deterius commutatum fuisse: anathema sit.}'' The canon is definitive teaching: a conciliar canon with an attached anathema.}

The weight is unambiguous---this is definitive teaching, a canon with an anathema---and so is the content. Adam \emph{incurred} death; the incurring was consequent on the offence; the death in question is the one previously threatened, which sends the reader back to Genesis 2:17; and the whole Adam, body and soul, was changed for the worse. Canon 2 adds the transmission, and adds it against a specific error: it anathematizes the man who says that Adam ``has only transfused death, and pains of the body, into the whole human race, but not sin also, which is the death of the soul.''\endnote{Trent, session V, canon 2, Waterworth English and Tauchnitz Latin, read as above. The canon's shape is instructive: the error it condemns is not the claim that bodily death was transmitted, which it takes for granted, but the claim that \emph{only} bodily death was. Its scriptural warrant is Romans 5:12 in the Vulgate form.} That canon takes bodily transmission as common ground and defines the further point.

Now notice what the canons do not say. They do not say that Adam's body was of itself incorruptible. They do not say how the exemption worked. They do not say whether an unfallen humanity would have been translated, transformed, or simply preserved. They do not adjudicate between the theological accounts of the preternatural gifts. Trent defined the \emph{that}, at the level of the sinner's liability, and left the \emph{how} to the schools---which is exactly what a council does when the error before it concerns liability and not physics.

That reticence is the space in which the Catechism's apparently contradictory sentences live.

\subsection{The distinction the schools worked out}

Aquinas takes the problem twice, and the second treatment is the one that has shaped Catholic usage ever since.

He asks first whether death and the other bodily defects are the \emph{effect} of sin, and gives an answer whose precision is easy to miss. Sin is not the cause of death \latin{per se}---no sinner intends to become mortal, and an effect not intended by the agent is not produced by the agent's own power. Sin is the cause of death \emph{accidentally}, in the specific sense of an agent that causes by removing an obstacle: ``it is stated in \emph{Phys.}~viii\ldots{} that `by displacing a pillar a man moves accidentally the stone resting thereon.' In this way the sin of our first parent is the cause of death and all such like defects in human nature, in so far as by the sin of our first parent original justice was taken away.''\endnote{ST I-II, q.~85, a.~5, corpus; English from the English Dominican translation, read 2026-07-26 at New Advent. The \latin{sed contra} is Romans 5:12. The reply to the first objection extends the image and is worth having: ``when original justice is removed, the nature of the human body is left to itself, so that according to diverse natural temperaments, some men's bodies are subject to more defects, some to fewer, although original sin is equal in all.'' The reply to the second explains why the baptized still die: our bodies remain subject to suffering ``in order that we may merit the impassibility of glory, in conformity with Christ.''}

The stone was always going to fall. It was resting on something. The sin knocked out what it was resting on.

Then the harder article: are death and such defects \emph{natural} to man? Aquinas stacks the objections on the side of nature---corruptible and incorruptible differ in genus, man is in the same genus as other animals, a body composed of contraries has the cause of its dissolution inside it---and the contrary arguments on the side of the form, including the flat scriptural assertion that ``God made not death.'' The resolution divides the nature rather than the answer: ``as regards his form, incorruption is more natural to man than to other corruptible things. But since that very form has a matter composed of contraries, from the inclination of that matter there results corruptibility in the whole. In this respect man is naturally corruptible as regards the nature of his matter left to itself, but not as regards the nature of his form.''\endnote{ST I-II, q.~85, a.~6, corpus, English from the English Dominican translation, read 2026-07-26 at New Advent; Latin verified the same day in the Corpus Thomisticum: ``\latin{Unde ex parte suae formae, naturalior est homini incorruptio quam aliis rebus corruptibilibus. Sed quia et ipsa habet materiam ex contrariis compositam, ex inclinatione materiae sequitur corruptibilitas in toto. Et secundum hoc, homo est naturaliter corruptibilis secundum naturam materiae sibi relictae, sed non secundum naturam formae.}'' The article's structure is unusual and deliberate: the three objections argue from matter, the three contrary arguments from form, and the corpus grants both.}

Then the craftsman. A smith choosing iron for a knife chooses hardness and ductility; that the iron also rusts ``results from the natural disposition of iron, nor does the workman choose this in the iron, indeed he would do without it if he could.'' So with the human body: nature chose a temperate mixture so the body could be an organ of touch; ``that it is corruptible is due to a condition of matter, and is not chosen by nature: indeed nature would choose an incorruptible matter if it could. But God, to Whom every nature is subject, in forming man supplied the defect of nature, and by the gift of original justice, gave the body a certain incorruptibility.\ldots{} It is in this sense that it is said that `God made not death,' and that death is the punishment of sin.''\endnote{ST I-II, q.~85, a.~6, corpus (continuation); Latin verified 2026-07-26 at the Corpus Thomisticum: ``\latin{Sed Deus, cui subiacet omnis natura, in ipsa institutione hominis supplevit defectum naturae, et dono iustitiae originalis dedit corpori incorruptibilitatem quandam\ldots{} Et secundum hoc dicitur quod Deus mortem non fecit, et quod mors est poena peccati.}'' The knife-and-rust comparison is Aquinas's, not this essay's; the same comparison recurs with a saw at ST II-II, q.~164, a.~1, ad~1.} And, in the parallel treatment, the whole thing compressed into one clause: ``death is both natural on account of a condition attaching to matter, and penal on account of the loss of the Divine favor preserving man from death.''\endnote{ST II-II, q.~164, a.~1, ad~1, English Dominican translation, read 2026-07-26 at New Advent. The reply expressly cross-references I-II, q.~85, a.~6. It should be received as school theology of high authority, not as defined doctrine: no council has canonized the matter/form solution, and the distinction between what is owed to a nature and what is superadded to it is a scholastic instrument, not an article of faith.}

This is common theology, and it should be labeled as such. The magisterium has not defined the matter/form solution; it has defined the liability, and then---in the Catechism, in \emph{Gaudium et spes}---spoken in language that the solution makes intelligible. \emph{Gaudium et spes} 18 says ``that bodily death from which man would have been immune had he not sinned''; the Catechism says ``man's nature is mortal'' and ``God had destined him not to die.'' A reader who has the distinction in hand finds those sentences unremarkable. A reader who does not finds them incoherent, and reasonably.

\begin{studybox}{Project synthesis: two questions, not two physics}
This is the essay's own judgment, argued from the sources cited above and attributed to no authority. The apparent contradiction between ``death is natural'' and ``death is a punishment'' dissolves once one sees that the two statements answer different questions, and only different questions. \emph{What kind of thing is this event?} is a question about the constitution of the thing that undergoes it, and the answer is: the dissolution of a composite whose matter is composed of contraries---the same answer one would give for an oak or an ox. \emph{Why is it happening to you?} is a question about the history of a relationship, and the answer is: because a gift that held the dissolution off is not being given. Neither answer competes with the other, because neither is in the other's business. The confusion arises entirely from hearing the second answer as a rival account of the first---as though ``death is a punishment'' were a claim about biology, and the Christian were saying that human bodies would not decompose if not for a transgression in a garden.

\medskip\noindent\textbf{Three reasons for the reading.} First, the tradition's own texts index the two answers to different principles \emph{of one nature}: the corruptibility to the matter, the incorruption to the form. That is a distinction inside a single analysis, not two competing histories of the world. Second, the magisterium defines only one side. Trent defines that Adam incurred the death that had been threatened; it says nothing about the constitution of his body, and this silence is not an oversight in a document that is elsewhere extremely specific. Third, the grammar of the sources is the grammar of gift: \latin{supplevit defectum naturae}, God \emph{supplied the defect of nature}. One does not supply what was already there, and the withdrawal of a supplement does not alter the nature it was supplementing. The technical name for what was removed---original justice, a gift---is itself the argument.

\medskip\noindent\textbf{The strongest objection, at full strength.} The distinction is a verbal maneuver that buys consistency by relocating the contradiction. Grant the Thomistic solution and ``natural'' has come to mean \emph{following from the matter considered in abstraction from the actual divine institution}---that is, from a nature that no human being has ever had. Adam as actually constituted was not mortal; his descendants are; the difference is a real difference in the world, brought about by a sin. To call the resulting condition ``natural'' is then to describe the actual human situation by reference to a counterfactual human being who never existed, and to let that fiction do the reconciling work. Worse, the objection has Scripture on its side at the one place where Scripture speaks directly to the point: \latin{Deus mortem non fecit}---God did not \emph{make} death. If bodily death follows from the matter God chose, then God made it, and the verse is false or means something other than what Aquinas takes it to mean. The synthesis above, on this objection, is not a resolution but an agreement to change the subject whenever the pressure comes from the other side.

\medskip\noindent\textbf{What would change the judgment.} A magisterial text, or an unambiguous locus in Aquinas, treating bodily immortality as \emph{owed} to human nature---\latin{debitum naturae} rather than \latin{donum}---would collapse the distinction and make the two accounts genuine rivals, because a debt is not a supplement. So would a demonstration that Trent's first canon defines not merely that Adam incurred death but that his body was of itself incorruptible; the essay has read the canon in both languages and finds no such definition, but the reading is of the promulgated text alone and not of the conciliar acts, and the acts might show the fathers deliberately intending more than the canon says. Either finding would require the synthesis to be withdrawn rather than qualified.
\end{studybox}

\subsection{What the distinction costs}

Nothing in the foregoing makes death less serious, and it is worth being explicit about that, because a distinction that dissolves a contradiction can look like a distinction that dissolves a doctrine.

It does not. On the reading defended here, every element of the Church's teaching survives intact: that man was made for a destiny incompatible with dissolution, that he does not now attain it by the road he was meant to travel, that the diversion is his own doing, and that what he now undergoes is a real deprivation and not a neutral biological terminus. Aquinas's own pillar-and-stone image makes the point sharply: it is no comfort to the man under the stone that gravity was always there. The gift was the point. That the gift's absence rather than its presence is what needs explaining is precisely the Christian claim, and it is a strange claim, and it should not be domesticated by the distinction that renders it consistent.

What the distinction does buy is the ability to say something honest to a person who is watching a body fail. It is not necessary to tell him that decay is unnatural; he can see that it is exactly what bodies do. It is necessary to tell him that this is not the arrangement he was made for, and that the difference between what he is watching and what should have been is the measure of a loss that is being addressed by other means.

Those other means are the second half of the Christian claim about death, and they change what death \emph{means} without changing a single one of its facts.
